\title{LongRoPE: Extending LLM Context Window Beyond 2 Million Tokens}

\begin{abstract}
A large context window is a sought-after characteristic in large language models (LLMs). Nonetheless, extending context length faces significant challenges, including prohibitive fine-tuning expenses, limited availability of lengthy text data, and the emergence of catastrophic values from novel token positions. Consequently, the practical context window size of current LLMs remains capped at approximately 128k tokens.

In this study, we present LongRoPE, which, for the first time, enlarges the context window of pre-trained LLMs to an impressive \textbf{2048k} tokens. Remarkably, this extension is accomplished using as few as 1k fine-tuning steps within training lengths of up to 256k, all while preserving the model’s original short-context performance. LongRoPE achieves these advances through three main contributions: (i) we identify and leverage two distinct types of non-uniformities in positional interpolation by employing an efficient search strategy. This yields superior fine-tuning initialization and permits an 8$\times$ window extension even without fine-tuning; (ii) we propose a progressive expansion method, wherein we first fine-tune an LLM on 256k token sequences, followed by a secondary positional interpolation applied to the extended, fine-tuned LLM to ultimately support 2048k token contexts; (iii) we re-tune LongRoPE at an 8k length to restore and uphold short-context performance. Comprehensive empirical studies on LLaMA2 and Mistral across diverse benchmarks confirm the efficacy of our approach. With only minimal adjustments to the positional embedding, models modified via LongRoPE retain their original architecture and are compatible with most pre-existing optimizations. The code is scheduled for release at \texttt{https://github.com/microsoft/LongRoPE}
\end{abstract}

\section{Introduction}

Although Large Language Models (LLMs) have achieved impressive outcomes across a range of tasks~\cite{gpt4,llama2}, they are frequently constrained by a restricted context window; for instance, LLaMA2 supports a maximum of 4096 tokens~\cite{llama2}. When input sequences exceed this limit, the effectiveness of LLMs diminishes, as these extended positions fall outside the models' training distribution. This limitation presents difficulties in key applications such as in-context learning where numerous examples are provided~\cite{Huang2023FewerIM}, and in the operation of LLM agents~\cite{park2023generative,madaan2023selfrefine}.

Recent studies demonstrate that the context window of a pre-trained LLM can be pushed to approximately 128k tokens via fine-tuning with lengthier sequences~\cite{longlora,pi,yarn,zhang2024soaring,liu2023scaling}. Nonetheless, three principal challenges limit further expansion of the context window. Firstly, newly introduced position indices, which are initially untrained, often bring many abnormal values, causing distribution shifts and complicating the fine-tuning process~\cite{pi}. The difficulty is amplified when expanding the window from 4k to $>$1000k, as this entails adding more than 90\% fresh positional indices. Secondly, fine-tuning generally necessitates texts that match the length of the desired window. However, current datasets offer limited supply of ultra-long texts, particularly those surpassing 1000k tokens. Additionally, utilizing such lengthy texts in training is computationally demanding, requiring substantial time and considerable GPU resources. Thirdly, as context windows become extremely large, attention mechanisms struggle to allocate focus effectively across the extensive range of token positions, thereby diminishing performance on tasks involving shorter contexts~\cite{pi}.

To address the initial challenge, one commonly used strategy is to interpolate RoPE positional embeddings~\cite{rope,pi}, whereby new position indices are rescaled to fit within the pretrained interval, as illustrated in Fig.\ref{fig:overview}. Position Interpolation (PI)~\cite{pi} executes linear interpolation of RoPE's rotary angles based on the desired extension ratio. The NTK method~\cite{ntk,dynamicntk} proposes an unequal scheme for both interpolation and extrapolation along different RoPE dimensions. YaRN~\cite{yarn} partitions RoPE dimensions according to three distinct frequency ranges, and then applies separate extrapolation, NTK-style, and linear interpolation approaches to each group. Despite these methods, positional embeddings in the Transformer framework are characterized by \emph{intricate non-uniform information entropy}. This nuanced heterogeneity is not adequately exploited by current interpolation techniques, resulting in information degradation and constraining the potential length of the context window.

Section~\ref{sec:analysis} presents two primary empirical insights: \textbf{(1)} Successful positional interpolation must account for two types of non-uniformity—specifically, the differences across RoPE dimensions as well as across token positions. Less interpolation is advantageous for lower RoPE dimensions and earlier tokens, although the best strategy is contingent on the intended length of extension. \textbf{(2)} Incorporating these non-uniform characteristics during positional interpolation preserves the informational integrity of the original RoPE—particularly across the most significant dimensions and token locations. As a result, this strategy reduces interpolation-related information loss and yields superior initializations for downstream fine-tuning. Furthermore, this approach supports up to an 8$\times$ increase in sequence length even in the absence of fine-tuning.

Building on these insights, we present \textbf{LongRoPE}, a robust approach that pushes the LLM context window size past 2 \emph{million} tokens. The methodology behind LongRoPE incorporates three main advances. The first is that {LongRoPE} leverages multidimensional, non-uniform positional interpolation in its entirety. It systematically determines suitable rescale factors for each RoPE dimension’s rotation angles, depending on the position of tokens. Given that the search space for these rescale factors grows exponentially with increased extension ratios, {LongRoPE} adopts an evolutionary search algorithm complemented by two optimization strategies to enhance the efficiency of this search. Figure~\ref{fig:overview} provides an illustration of the resulting rescaled RoPE derived from this process.

Next, {LongRoPE} adopts an efficient, stepwise extension approach that enables expanding the context window to 2048k, circumventing the requirement to fine-tune directly on extremely lengthy texts, which are both uncommon and difficult to obtain. The method initiates by identifying a suitable rescale factor for a 256k sequence length on the pre-trained LLM and performing fine-tuning with this sequence length. Subsequently, thanks to the capacity of our non-uniform positional interpolation to support up to an 8$\times$ increase in non-fine-tuned scenarios, a second stage search is carried out to determine additional RoPE rescale factors using the already fine-tuned, length-extended LLM. This approach successfully extends the context window to 2048k for LLaMA2 and Mistral~\cite{mistral}.

In order to prevent diminished performance when operating within the original, smaller context window, {LongRoPE} maintains the adjustment of RoPE rescale factors in the expanded LLM. Following the approach used for increasing context length from 256k to 2048k, we apply our search algorithm to decrease the context window to 4k and 8k on the LLM fine-tuned with 256k, thereby reducing positional interpolation. For inference, when the sequence length falls below 8k, RoPE is modified with the optimized rescale factors discovered through this search.

Through comprehensive testing on multiple large language models and diverse tasks requiring extended context, the efficacy of our approach is well established. Our findings indicate that LongRoPE consistently preserves low perplexity across evaluation lengths ranging from 4k to 2048k, surpasses 90\% in passkey retrieval accuracy, and yields performance equivalent to conventional benchmarks using a 4096-token context window. LongRoPE is compatible with all LLMs employing RoPE embeddings. We plan to make our code and LongRoPE-2048k models publicly available.

\section{Non-uniformity in Positional Interpolation}
\label{sec:analysis}

\subsection{Preliminary}

Transformer models require explicit positional information, often in the form of position embedding, to represent the order of input tokens. Our work focuses on the RoPE~\cite{rope} position embedding, which is widely used in recent  LLMs. For a token at position index $n$, its corresponding RoPE encoding can be simplified as follows:

\begin{equation}
		\fontsize{7.8}{7.8} \selectfont
	\label{eq:rope}
	\left[ cos(n\theta_0), sin(n\theta_0), cos(n\theta_1), \cdot\cdot\cdot, cos(n\theta_{d/2-1}), sin(n\theta_{d/2-1}) \right]   
\end{equation}

Here, $d$ denotes the dimension of embeddings, 
$n\theta_i$ refers to the rotary position angle for the $n$-th token, 
and $\theta_i=\theta^{-2i/d}$ specifies the rotational frequency. In RoPE, $\theta$ typically defaults to 10000.

\textbf{\emph{Context window extension ratio $s$} and positional interpolation}. The parameter $s$ is defined as the quotient of the new context length $L'$ over the initial length $L$, i.e., $s=\frac{L'}{L}$.

To extend the context window from $L$ to $L'$, current positional interpolation methods suggest downscaling rotation frequencies $\theta_i$ based on extension ratio $s$. Let $\beta=\theta^{2/d}$, and $\lambda$ denote the actual rescale factor related to $s$, we   unify these positional interpolation methods as follows:

\begin{equation}	
	\fontsize{7.5}{7.5} \selectfont
	\label{eq:unified_rope}
	\begin{aligned}
		\left[cos\left(\frac{n}{\lambda(\beta)^0}\right), sin \left(\frac{n}{\lambda(\beta)^0}\right), cos\left(\frac{n}{\lambda(\beta)^1}\right), \cdot\cdot\cdot,  sin \left(\frac{n}{\lambda(\beta)^{d/2-1}}\right) \right]   
	\end{aligned}
\end{equation}

\textbf{Linear positional interpolation (PI)}. PI~\cite{pi} involves interpolating position indices in a linear fashion, staying within the original length used for pre-training. For a given extension factor $s$, this method uniformly scales down the rotation angles at every position by $\lambda=s$ throughout all RoPE dimensions. Nevertheless, this linear compression leads to positional representations being densely packed together, making it difficult for the model to discriminate between tokens with nearby positions. As a consequence, PI's effectiveness typically diminishes when the extension ratio becomes large.

\textbf{NTK-based interpolation and extrapolation}. The works of ~\cite{ntk,dynamicntk} investigate RoPE using an information encoding lens and leverage Neural Tangent Kernel (NTK) theory~\cite{jacot2018neural,tancik2020fourier}. To address the problem of overcrowded positions inherent in PI, these studies propose distributing the interpolation burden more evenly among the RoPE dimensions. By reducing the scaling of the lower (high-frequency) dimensions and increasing the scaling of the higher (low-frequency) ones, this approach enables both positional interpolation and extrapolation, characterized by $\lambda=s^{i}$. The dynamic NTK extension~\cite{dynamicntk} further refines this method by dynamically modifying the extension ratio for each position according to the current sequence length. Unlike PI, which typically requires additional fine-tuning, NTK-based techniques offer the ability to expand context windows even without fine-tuning, though they generally cap the maximum extension ratio at 4$\times$.

\textbf{YaRN}~\cite{yarn} presents notable advancements in positional interpolation by categorizing RoPE dimensions into three distinct frequency-based segments, each assigned a unique interpolation method. Extrapolation ($\lambda$=1) is utilized for dimensions with high frequencies, linear interpolation (PI) is applied to those with low frequencies, and the NTK approach is employed for intermediate frequencies. The central innovation in YaRN is this dimensional grouping strategy; however, it is currently based on empirical determination by researchers. This reliance on manual selection could limit YaRN's effectiveness for novel LLM architectures.

\subsection{Study on Non-uniform Positional Interpolation}

Drawing on insights from NTK and YaRN, we observe improvements derived from incorporating non-linearity, particularly through the treatment of distinct frequencies along RoPE dimensions that facilitate targeted interpolation and extrapolation. Nevertheless, prevailing non-linear approaches are predominantly governed by manually crafted heuristics. Consequently, this prompts two pertinent inquiries: (1) Does the existing method for positional interpolation represent the best possible solution? (2) Might there be alternative non-linear strategies that have yet to be explored?

To investigate these questions, we apply evolution search (refer to Sec.~\ref{sec:method}) to identify improved non-uniform positional interpolations for LLaMA2-7B. The search process is steered by perplexity results, employing 5 randomly selected instances from the PG19~\cite{pg19} validation dataset. Our experiments lead us to the following important observations.

\textbf{Finding 1}: \textit{RoPE dimensions exhibit substantial non-uniformities, which are not effectively handled by current positional interpolation methods.}

We determine the optimal $\lambda$ for each RoPE dimension specified in Eq.~\ref{eq:unified_rope}. Table~\ref{tbl:exp1} presents the perplexity results of LLaMA2-7B evaluated with various approaches on the PG19 and Proof-pile~\cite{proof-pile} test sets, all without any fine-tuning. Our optimization method yields noticeable gains, indicating that both linear (PI) and non-uniform (Dynamic-NTK and YaRN) interpolation strategies currently fall short of optimality. Interestingly, YaRN performs worse than both PI and NTK on PG19, attributed to its failure to achieve the desired context window length in LLMs that have not been fine-tuned. As an illustration, with an 8k context size, YaRN’s perplexity sharply increases after reaching 7k.

Our investigation reveals that the rescaling coefficients $\lambda$ in Eq.~\ref{eq:unified_rope} are not constant, in contrast to the uniform scaling factor $s$ employed in PI, NTK’s analytical methods, and YaRN’s group-based computations. The adoption of these variable scaling factors yields substantial improvements in LLaMA2’s language modeling ability (measured by perplexity) across extended context lengths of 8k and 16k, without necessitating fine-tuning. This enhancement is attributed to the modified positional embedding’s capacity to maintain the essential characteristics of the original RoPE—particularly in crucial dimensions—thereby facilitating the model’s discrimination between tokens situated at proximate positions.

\textbf{Finding 2}: \textit{RoPE for the initial tokens in the input sequence should be extrapolated with less interpolation.} 

We propose that for the first $\hat{n}$ tokens in input sequences, the RoPE mechanism should apply reduced interpolation. The rationale behind this is that these tokens typically attract high attention scores, signifying their importance within attention layers—a trend previously noted in Streaming LLM~\cite{streamingllm} and LM-Infinite~\cite{han2023lminfinite}. To empirically test this assumption, we expand the context window to 8k and 16k via PI and NTK methods, systematically excluding the first $\hat n$ tokens (where $\hat n$=0,2,...,256) from any interpolation. The setup with $\hat n$=0 corresponds to the conventional PI and NTK. As summarized in Table~\ref{tbl:tokenposition}, our experiments yield two principal findings: \textbf{(1)} preserving the initial tokens without positional interpolation leads to a measurable performance gain; \textbf{(2)} the optimal value for $\hat n$ varies according to the specific extended context length desired.

\textbf{Finding 3}: \textit{Non-uniform positional interpolation effectively extends LLM context window in both fine-tuning and non-fine-tuning settings.} 

Our findings indicate that applying our optimized non-uniform position interpolation yields notable gains in extension capabilities at context lengths of 8k and 16k, even in the absence of fine-tuning. However, for extending to longer context windows, model fine-tuning becomes necessary. Consequently, we perform fine-tuning on LLaMA2-7B employing the searched RoPE configuration for a 64k context window (refer to Appendix for experimental details). As presented in Table~\ref{tbl:finetune}, our strategy consistently surpasses both PI and YaRN approaches, whether prior to or following the fine-tuning of LLaMA2-7B. The superior results stem from our method’s use of non-uniform positional interpolation, which helps minimize information degradation and delivers a more advantageous starting point for the fine-tuning process.

\textbf{Summary}. Through our research, we reveal two types of non-uniformities—across RoPE dimensions and among token positions. Harnessing these non-uniformities in positional interpolation proves highly effective in substantially enhancing LLM performance for extended context lengths.

\section{LongRoPE}

\label{sec:method}
Drawing on these observations, we propose LongRoPE, a method that begins by implementing an optimized search strategy to effectively leverage both forms of non-uniformity. This approach is subsequently utilized to expand the LLM context window to over 2 million tokens.

\subsection{Problem Formulation}

These two sources of non-uniformity greatly expand the solution space and add layers of difficulty to the optimization process. To mitigate this challenge, we reformulate the optimization task for multidimensional non-uniform position interpolation as a search problem.

For a LLM targeting a  context window size of $L'$ and lengthy input documents $\mathbf{X}$, where each $\mathbf{x}\in\mathbf{X}$ surpasses $L'$ in token length, we denote the original rotary angle of the $i^{th}$ dimension in RoPE embedding at token position $n$ as  $\frac{n}{\beta^i}$. The optimization problem is then formulated as follows:

\begin{equation}
\begin{gathered}
		\label{eq:problem}

\underset {\mathbf{x}\in\mathbf{X}; \, |\mathbf{x}|\ge L'}{ \operatorname {arg\,min} } \,  \mathcal{L}\, \left( \text{LLM} (\text{RoPE}, \mathbf{X})\right), \quad \text{where}
\\
\scalebox{0.93}{
$\underset {\begin{subarray}{l} i=0,\ldots,\frac{d}{2}-1;\\ n\in[0,|\mathbf{x}|);\end{subarray}}{\text{RoPE}_ {(n)}} = {\left[\ldots, \cos\left(\mathbb{I}(\hat{\lambda}_{i}, \hat{n})\cdot\frac{n}{\beta^{i}}\right), \sin\left(\mathbb{I}(\hat{\lambda}_{i}, \hat{n})\cdot\frac{n}{\beta^{i}}\right), \ldots\right]}$}
\\
\text{where} \quad \mathbb{I}(\hat{\lambda}_{i},\hat{n})=
\begin{cases}
1&\;\;\text{if}\; n<\hat{n}\\
\frac{1}{\lambda_{i}} &\;\;\text{if}\; n\geq\hat{n}
\end{cases}
\end{gathered}
\end{equation}

We define a set of rescale factors, $\mathbb{I}(\hat{\lambda}_{i},\hat{n})$, to address both types of non-uniformities present. Here, $\hat{\lambda_i}$ and $\hat{n}$ represent the dimension-related non-uniformity in RoPE and the positional non-uniformity among tokens, respectively. More concretely, the function $\mathbb{I}(\hat{\lambda}_{i},\hat{n})$ is employed to adjust the rotation angle in the $i^{th}$ RoPE dimension; $\hat\lambda_{i}$ is used as the scaling parameter, while $\hat{n}$ serves as the threshold for token positions. When the token position is less than $\hat{n}$, that is, for the first $\hat{n}$-1 tokens, $\hat\lambda_{i}$ is not applied, so the standard RoPE rotary angle, $\frac{n}{\beta^i}$, remains. When $n\ge\hat{n}$, the rescale factor $\hat\lambda_{i}$ is incorporated.

Suppose we set a desired context window length of $L'$. The aim is then to identify the most suitable rescale factors ($\mathbb{I}(\hat{\lambda}_{0},\hat{n})$, $\mathbb{I}(\hat{\lambda}_{1},\hat{n})$, ..., $\mathbb{I}(\hat{\lambda}_{i},\hat{n})$, ...) corresponding to each of the RoPE dimensions from $1$ to $d$. By appropriately rescaling $\text{RoPE}$ within the target $\text{LLM}$, we seek to minimize the next token prediction loss $\mathcal{L}$ (which also corresponds to perplexity) for input sequences $\mathbf{X}$ consisting of $L'$ tokens.

\subsection{Searching the Non-uniform Position Interpolation}

In addressing the challenge presented in Eq.~\ref{eq:problem}, we propose a straightforward but powerful approach that identifies the ideal RoPE rescale factors. This enables comprehensive utilization of the varying multidimensional characteristics found within position embeddings.

\textbf{Search space}. We construct an expansive search space that accommodates both types of non-uniformity. The structure of this search space is depicted in Table~\ref{tbl:searchspace}. In detail, we permit the optimization of individual rescale factors for each RoPE embedding dimension. To streamline the process, the search is performed over $\lambda_i$ and $\hat{n}$, rather than directly exploring $\mathbb{I}(\hat{\lambda}_{i},\hat{n})$, with $\hat{\lambda}_{i}$ defined as $1/\lambda_i$. According to Table~\ref{tbl:searchspace}, the search for each $\lambda_i$ ranges from a lower bound of 1.0 (representing pure extrapolation) up to a maximum of $s\times1.25$ (indicating more pronounced interpolation than PI), with increments set at 0.01. Here, $s$ denotes the proportion by which the context window is extended.

$\hat{n}$ determines how many of the beginning token positions are preserved with their original positional encodings, bypassing position interpolation (specifically maintaining the native RoPE embedding). In practice, we conduct a search for $\hat{n}$ over the set \{0, 1, 2, 4, 8, 12, 16, 20, 24, 28, 32, 64, 128, 256\}. If $\hat{n}=0$, then the rescale factors found through the search are applied to all token positions.

\textbf{Evolution-based search}. The search space detailed in Table~\ref{tbl:searchspace} encompasses a vast array of positional interpolation configurations, creating substantial difficulty in achieving efficient exploration. For instance, increasing the extension ratio to $s=4\times$ results in $400^{128/2}\times14=4\times10^{167}$ possible combinations. As the extension ratio grows, this space increases at an exponential rate. To manage this complexity, we employ evolution search~\cite{guo2020single}, and further implement two optimization strategies that markedly enhance search efficiency. The overall search methodology is presented in Algorithm~\ref{alg:search}.

\textit{Enhanced initialization of population}. Rather than randomly assigning all $P$ rescale factors at the outset, we explicitly include the three RoPE rescale factors tied to PI, NTK, and YaRN among the initial members. For the other $P$-3 individuals, these three rescale factors are subjected to random mutation with probability $p$.

\textit{Monotonically non-decreasing constraint}. Once the initial population is formed, we evaluate each candidate by measuring LLM perplexity. This is done by applying their specific RoPE rescale parameters to the target LLM and then calculating perplexity on input $\mathbf{X}$. The highest-ranked k individuals are chosen as progenitors for subsequent generations. Nevertheless, the expansive search landscape means that basic mutation and crossover operations often yield suboptimal candidates, which results in superfluous perplexity evaluations. This inefficiency becomes more pronounced as $L'$ increases, since every perplexity assessment requires costly inference.

To tackle this issue, we enforce a monotonic non-decreasing constraint on the rescaled RoPE factors sampled, specifically ensuring $\lambda_i\le \lambda_{i+1}$. RoPE configurations that adhere to this constraint are the only ones utilized in LLM perplexity assessments, thereby greatly curtailing the computational demands of searching. In particular, the constraint stipulates that $\lambda_i$ must exhibit monotonically increasing values as the RoPE dimension grows ($i = 0,\ldots,63$). This pattern of dimensional monotonicity is informed by NTK theory~\cite{jacot2018neural,tancik2020fourier,ntk}, which posits that lower-frequency dimensions, associated with higher frequencies, benefit from minimal interpolation (i.e., smaller $\lambda_i$), whereas higher-frequency dimensions—corresponding to lower frequencies—permit greater interpolation (i.e., larger $\lambda_i$).

\begin{algorithm}[t]
	\small
	\caption{The search algorithm}
	\textbf{Input:} target LLM, input samples $\mathbf{X}$, population size $P$, mutation size $N_1$, crossover size $N_2$, maximum iterations $\mathcal{T}$, mutation probability $p$ \\

	\begin{algorithmic}[1]
		\label{alg:search}
		\STATE Top-k $\leftarrow \phi$;	
		\STATE $\text{P}_0 \leftarrow$ \textit{Initialize\_population\_with\_optimization}($P$, $\mathbf{X}$, $p$); 
		\FOR{$i=1$ to $\mathcal{T}$}
		\STATE \textit{Compute\_perplexity}(LLM, $\text{P}_{i-1}$, $\mathbf{X}$);
		\STATE Top-k $\leftarrow$ \textit{Update\_Topk}(Top-k, $\text{P}_{i-1}$);
		\STATE $\text{P}_{mutation} \leftarrow$ \textit{Mutation\_with\_mono\_constraint}(Top-k, $N_1$, $p$);
		\STATE $\text{P}_{crossover} \leftarrow$ \textit{Crossover\_with\_mono\_constraint}(Top-k, $N_2$);
		\STATE $\text{P}_i \leftarrow \text{P}_{mutation} \cup \text{P}_{crossover} \cup$ Top-k;
		\ENDFOR
		\STATE Return the member in Top-k with minimal perplexity;
	\end{algorithmic}
\end{algorithm}

\textbf{8$\times$ extension achieved without fine-tuning}. Through the use of an evolutionary search strategy, our approach selects optimal non-uniform RoPE rescale factors, ensuring the retention of essential dimensions and positions to reduce information loss from interpolation. As illustrated in Fig.\ref{fig:non-ft}, this technique enables LLaMA2's context window to expand from 4k to 32k without necessitating fine-tuning. Conversely, prior approaches including PI, non-uniform NTK, and YaRN experience a sharp increase in perplexity upon merely doubling the context window.

\subsection{Extending LLM Context Window to 2048K}
\label{sec:extendto2048k}

\textbf{Progressive extension to 2048k}. In this section, we present our approach for expanding the context window of pre-trained LLMs from the standard 4k length up to more than 2048k tokens. Our non-uniform positional interpolation technique is capable of delivering up to an 8$\times$ increase in context length without necessitating fine-tuning. However, achieving significantly larger extensions (for example, 512$\times$) does require fine-tuning. A potential strategy involves searching for suitable RoPE rescale factors corresponding to the target context window of 2048k, followed by fine-tuning. Nevertheless, this approach entails extremely high computational and training costs. In addition, based on our experience, achieving effective fine-tuning at such large extension ratios presents considerable difficulty (refer to the Appendix for further details).

Fortunately, LongRoPE demonstrates efficacy with both the base and fine-tuned, extended LLMs. To this end, we propose an efficient, stepwise approach capable of reaching the intended 2048k sequence length using only 1k fine-tuning steps, all conducted at a maximum training length of 256k.

$\Diamond$ \textit{LongRoPE search for scaling pre-trained LLM to 256k}. Using LLaMA2 as a case study, we perform searches aimed at expanding the context window to 128k and 256k tokens. At this phase, the model undergoes extension by factors of 32$\times$ and 64$\times$, respectively.

$\Diamond$ \textit{Fine-tuning to 256k}. Subsequently, we adapt the pre-trained LLM to support a 256k context window. In detail, LLaMA2 undergoes an initial fine-tuning phase of 400 steps employing RoPE rescaled factors tailored to 128k. Once this stage concludes, we update the RoPE rescaled factors to match 256k on the resulting checkpoint and proceed with a further 600 fine-tuning steps. This two-step approach is found to be more effective than immediately fine-tuning for a 256k context.

$\Diamond$ \textit{Increasing fine-tuned extended LLM to 2048k using LongRoPE search}. In the last step, we carry out another search process with the LLM that was fine-tuned for sequences of length 256k. Through this approach, we successfully achieve a remarkable context window of 2048k, accomplishing this expansion without additional rounds of fine-tuning. The overall context window enlargement amounts to $512\times$.

\textbf{Recovery of Performance in Shorter Context Windows}. When scaling up the context window to 2048k, there is a noticeable decline in model effectiveness within the initial context span. This phenomenon has been observed previously in positional interpolation~\cite{pi}: projecting position embeddings into higher-dimensional spaces compresses their representation in the original window, which hampers language modeling accuracy. Specifically, with an extension ratio of 512$\times$, the spatial representation of positions in the primary 4k context window becomes excessively dense.

To address this issue, we conduct an additional evolutionary search on the extended LLM, optimizing RoPE rescale parameters specifically for shorter context lengths such as 4k and 8k. The upper bound for searched $\lambda$ is decreased since shorter contexts necessitate minimal positional interpolation. At inference time, the LLM updates the relevant RoPE rescale parameters dynamically.

\section{LongRoPE}

\label{sec:method}
Building on these insights, we propose LongRoPE, a method that initially implements an optimized search technique to take full advantage of the two forms of non-uniformity. This approach is subsequently employed to expand the LLM context window to accommodate over 2 million tokens.

\subsection{Problem Formulation}

These two forms of non-uniformity expand the solution space considerably and add layers of complexity to the optimization process. To manage this challenge, we recast the optimization of multidimensional non-uniform position interpolation as a search-based problem.

For a LLM targeting a  context window size of $L'$ and lengthy input documents $\mathbf{X}$, where each $\mathbf{x}\in\mathbf{X}$ surpasses $L'$ in token length, we denote the original rotary angle of the $i^{th}$ dimension in RoPE embedding at token position $n$ as  $\frac{n}{\beta^i}$. The optimization problem is then formulated as follows:

\begin{equation}
\begin{gathered}
		\label{eq:problem}

\underset{\mathbf{x} \in \mathbf{X};\, |\mathbf{x}| \geq L'}{\operatorname{arg\,min}}\, \mathcal{L} \left( \text{LLM}(\text{RoPE}, \mathbf{X}) \right),\quad\text{where}
\\
\scalebox{0.93}{
$\underset{\begin{subarray}{l} i = 0, \ldots, \frac{d}{2} - 1; \\ n \in [0, |\mathbf{x}|);\end{subarray}}{\text{RoPE}_n} = \left[ \ldots, \cos\left( \mathbb{I}(\hat{\lambda}_i, \hat{n}) \cdot \frac{n}{\beta^i} \right),\, \sin\left( \mathbb{I}(\hat{\lambda}_i, \hat{n}) \cdot \frac{n}{\beta^i} \right),\, \ldots \right]$
}
\\
\text{where}\quad \mathbb{I}(\hat{\lambda}_i, \hat{n}) = \begin{cases}
	1 &\quad n < \hat{n} \\
	{\frac{1}{\lambda}_i} &\quad n \geq \hat{n}
\end{cases}
	\end{gathered}
\end{equation}

In this framework, we define a collection of rescale coefficients, $\mathbb{I}(\hat{\lambda}_{i},\hat{n})$, to address both types of non-uniformity. Here, $\hat{\lambda_i}$ and $\hat{n}$ indicate the degree of non-uniformity present in the RoPE dimensions and the token positions, respectively. The purpose of $\mathbb{I}(\hat{\lambda}_{i},\hat{n})$ is to modify the rotation angle corresponding to the $i^{th}$ RoPE dimension, such that $\hat\lambda_{i}$ serves as the adjustment multiplier and $\hat{n}$ determines the threshold position. For the first $\hat{n}-1$ tokens, $\hat\lambda_{i}$ remains inactive, and the standard RoPE rotary angle, expressed as $\frac{n}{\beta^i}$, is maintained. When the token index reaches $n\ge\hat{n}$, the scaling factor $\hat\lambda_{i}$ is then engaged.

Assuming a desired context window length of $L'$, the task is to determine the best set of rescaling coefficients ($\mathbb{I}(\hat{\lambda}_{0},\hat{n})$, $\mathbb{I}(\hat{\lambda}_{1},\hat{n})$, ... $\mathbb{I}(\hat{\lambda}_{i},\hat{n})$, etc.) for each RoPE dimension, from the $1^{st}$ up to the $d^{th}$. This optimization enables the modified $\text{LLM}$, using the adjusted $\text{RoPE}$, to minimize the loss for next token prediction, denoted as $\mathcal{L}$ (which reflects perplexity), on datasets $\mathbf{X}$ that contain $L'$ tokens.

\subsection{Searching the Non-uniform Position Interpolation}

To address the issue outlined in Eq.~\ref{eq:problem}, we propose an efficient and straightforward approach that identifies optimal RoPE rescale factors, enabling the comprehensive utilization of multidimensional irregularities within position embeddings.

\textbf{Search space}. We construct an expansive search space to capture both forms of non-uniformity. Table~\ref{tbl:searchspace} provides an overview of this search space. In detail, we permit the optimization of individual rescale factors for each dimension in the RoPE embedding. To streamline the configuration, we focus the search on $\lambda_i$ and $\hat{n}$ rather than directly searching over $\mathbb{I}(\hat{\lambda}_{i},\hat{n})$, noting that $\hat{\lambda}_{i}=1/\lambda_i$. As presented in Table~\ref{tbl:searchspace}, the variable $\lambda_i$ can be selected from a lower bound of 1.0 (representing basic extrapolation) up to $s\times 1.25$ (corresponding to greater interpolation relative to PI), using increments of 0.01. Here, $s$ denotes the intended ratio for context window expansion.

The parameter $\hat{n}$ determines how many starting token positions are preserved in their original form, without applying position interpolation—specifically, maintaining the standard RoPE embedding. In practice, we let $\hat{n}$ be selected from the set \{0, 1, 2, 4, 8, 12, 16, 20, 24, 28, 32, 64, 128, 256\}. If $\hat{n}=0$, every token position adopts the learned rescale factors.

\textbf{Evolution-based search}. The search space outlined in Table~\ref{tbl:searchspace} encompasses a vast array of positional interpolation methods, making systematic exploration highly demanding. As an illustration, extending to $s=4\times$ results in $400^{128/2}\times14$, amounting to $4\times10^{167}$ possible configurations. When the extension ratio increases, the number of potential solutions grows exponentially. To efficiently traverse this extensive landscape, we leverage evolution search~\cite{guo2020single}, supplemented by two specific optimization strategies to markedly enhance search performance. The workflow of this approach is detailed in Algorithm~\ref{alg:search}.

\textit{Optimized initial population generation}. Rather than creating a population of $P$ rescale factors purely at random, we specifically include the three RoPE rescale factors associated with PI, NTK, and YaRN as members of the initial population. For the other $P$-3 individuals, we generate them by introducing random mutations to the three rescale factors, each mutation occurring with probability $p$.

\textit{Monotonically non-decreasing constraint}. Once the initial population is established, LLM perplexity is evaluated for each candidate. This involves adjusting the target LLM with the relevant RoPE rescale parameters and calculating the perplexity on the input $\mathbf{X}$. The highest-scoring k individuals are chosen as parents for further evolutionary steps. Due to the expansive nature of the search space, straightforward mutation and crossover can yield suboptimal solutions and result in excessive and redundant perplexity computations. This inefficiency becomes increasingly pronounced for larger values of $L'$, as every perplexity assessment demands substantial inference time.

To mitigate this, we enforce a non-decreasing monotonicity condition on the sampled RoPE rescaling factors, such that $\lambda_i \leq \lambda_{i+1}$. Only those RoPE configurations adhering to this restriction are utilized during perplexity assessments in the LLM, which substantially streamlines the search process. Concretely, we stipulate that $\lambda_i$ must monotonically increase across the RoPE dimensions ($i = 0, \ldots, 63$). This requirement is motivated by NTK theory~\cite{jacot2018neural,tancik2020fourier,ntk}, which implies that lower-indexed dimensions, corresponding to higher-frequency components, necessitate minimal interpolation (resulting in smaller values for $\lambda_i$), whereas higher-indexed dimensions, tied to lower-frequency components, can accommodate more interpolation (yielding larger $\lambda_i$).

\begin{algorithm}[t]
	\small
	\caption{The search algorithm}
	\textbf{Input:} target LLM, input samples $\mathbf{X}$,   population size $P$, mutation size $N_1$, crossover size $N_2$, max iterations $\mathcal{T}$, mutate probability $p$\\

	\begin{algorithmic}[1]
		\label{alg:search}
		\STATE Top-k$\leftarrow$ $\phi$;	
		\STATE $\text{P}_0$ = \textit{Initialize\_population\_with\_optimization} ($P$, $\mathbf{X}$, $p$); 
		\FOR{$i$ = $1$ to $\mathcal{T}$}
		\STATE \textit{Compute\_perplexity} (LLM, $\text{P}_{i-1}$, $\mathbf{X}$);
		\STATE Top-k $\leftarrow$ \textit{Update\_Topk} (Top-k, $\text{P}_{i-1}$);
		\STATE $\text{P}_{mutation}$ = \textit{Mutation\_with\_mono\_constraint} (Top-k, $N_1$, $p$); 
		\STATE $\text{P}_{crossover}$ = \textit{Crossover\_with\_mono\_constraint} (Top-k, $N_2$);
		\STATE $\text{P}_i$ = $\text{P}_{mutation}$ $\cup$ $\text{P}_{crossover}$ $\cup$ Top-k;
		\ENDFOR
		\STATE Output the member with minimum perplexity in Top-k;
	\end{algorithmic}
\end{algorithm}

\textbf{8$\times$ extension achieved without fine-tuning}. Through evolutionary search, our approach successfully selects non-uniform RoPE rescale factors, retaining critical dimensions and positional information to reduce information loss associated with interpolation. As shown in Fig.\ref{fig:non-ft}, this enables us to increase LLaMA2's context window from 4k to 32k without the need for fine-tuning. By comparison, previous techniques like PI as well as non-uniform NTK and YaRN lead to a sharp rise in perplexity beyond a 2$\times$ increase.

\subsection{Extending LLM Context Window to 2048K}
\label{sec:extendto2048k}

\textbf{Progressive extension to 2048k}. In this section, we present an approach for expanding the context window of pre-trained LLMs from the conventional 4k length up to more than 2048k tokens. Our results indicate that the proposed non-uniform positional interpolation can support an extension up to 8$\times$ the original window size without necessitating fine-tuning. For more substantial extensions, such as reaching 512$\times$, fine-tuning becomes indispensable. A typical strategy involves determining suitable RoPE rescale factors at the target length of 2048k, followed by a subsequent fine-tuning phase. Nonetheless, this approach encounters significant obstacles, primarily due to the extremely high computational demands. Additionally, as observed in our experiments, it remains difficult to effectively fine-tune LLMs when subjected to such drastic increases in context window size (refer to the Appendix for further details).

LongRoPE demonstrates efficacy with both the base and further fine-tuned extended LLM variants. Building on this, we propose a streamlined and incremental approach that enables attainment of a 2048k target context length, requiring only 1k fine-tuning iterations and operating within a training sequence of 256k tokens.

$\Diamond$ \textit{Scaling pre-trained LLMs to 256k utilizing LongRoPE search}. Using LLaMA2 as a reference model, we explore options for expanding the context window to 128k and 256k. At these points, the scaling factors are 32$\times$ and 64$\times$ respectively.

$\Diamond$ \textit{Fine-tuning to 256k}. Next, we adapt the pre-trained LLM to accommodate a 256k context window. Our approach involves an initial fine-tuning phase on LLaMA2 for 400 steps, utilizing RoPE rescaled factors set to 128k. Afterward, we update the RoPE rescaled factors to 256k in the completed checkpoint and continue with 600 more steps of fine-tuning. This staged process demonstrates greater efficiency compared to fine-tuning for 256k context length in one step.

$\Diamond$ \textit{Applying LongRoPE search to broaden fine-tuned extended LLM context to 2048k}. In the last step, we carry out an additional search on the already fine-tuned LLM set to a 256k context length. This process increases the model’s context window to a substantial 2048k, and notably, this is achieved without any additional fine-tuning. The overall context enlargement factor reaches $512\times$.

\textbf{Restoring performance in shorter context windows}. Upon expanding the context window to an extensive 2048k length, we observe a decline in model accuracy within the initial context window. This reduction is consistent with established challenges associated with positional interpolation~\cite{pi}. The technique compresses position embeddings corresponding to the original context window into a considerably restricted subset of the representational space, which diminishes the effectiveness of the language model. Specifically, under a 512$\times$ scaling factor, the positional embeddings within the starting 4k context window become densely clustered.

To address this issue, we conduct an additional evolutionary search on the extended LLM to fine-tune the RoPE rescale factors specifically for shorter context lengths (such as 4k and 8k). For these reduced lengths, the upper limit for the searched $\lambda$ is lowered, reflecting the decreased need for positional interpolation. When running inference, the LLM adaptively modifies the appropriate RoPE rescale factors as needed.

 \section{Experiments}

\label{sec:exp}
\subsection{Setup}
\textbf{Evaluation Tasks and models}. We implement LongRoPE with LLaMA2-7B and Mistral-7B models, assessing their effectiveness through three dimensions: (1) perplexity metrics for LLMs equipped for longer contexts when handling lengthy documents; (2) the Passkey retrieval challenge, designed to test the model's proficiency in extracting a straightforward passkey hidden amidst largely irrelevant passages; and (3) conventional LLM benchmarks conducted within the typical 4096-token context window.

\textbf{Fine-tuning}. In the case of LLaMA2, we employ a learning rate set at 2e-5 with a linear decay schedule and a global batch size of 32. The model undergoes fine-tuning for 400 steps using the Redpajama~\cite{redpajama} dataset, which is partitioned into 128k-token sequences bounded by BOS and EOS tokens. Building upon the obtained checkpoint, we continue training for an additional 600 steps to extend the context window to 256k tokens. Training with a 128k context size utilizes 8 A100 GPUs via the distributed training framework~\cite{cube}, while expanding to a 256k context requires 16 A100 GPUs. For Mistral, training is conducted with a fixed learning rate of 1e-6 and a global batch size of 64. Both the 128k and 256k variants follow the configurations described in YaRN~\cite{yarn}, involving 400 training steps on the Together Computer's Long-Data Collections~\cite{mistraldata} with a sequence length of 16k. The training process utilizes 4 A100 GPUs.

\textbf{Search}. When the target window size is up to 256k, we set parameters as follows: $P$=64, $N_1$=$N_2$=16, $p$=0.3, $\mathcal{T}$=40. During each iteration, the top-32 candidates are chosen for mutation and crossover steps. Perplexity evaluation is performed using 5 randomly chosen validation samples from PG19, ensuring each sample meets the minimum required context length. For cases with window sizes greater than 512k, the population as well as mutation and crossover counts are reduced by half. In such configurations, perplexity is assessed based on 3 random examples drawn from the Pile-Books3~\cite{pile} validation set.

\textbf{Baselines}. To achieve a context window of 2048k, models were fine-tuned initially with 128k and 256k context lengths. As a result, LongRoPE-2048k (ft=128k) was developed for LLaMA2, while LongRoPE-2048k (ft=256k) was produced for Mistral. These four models are evaluated against prominent context window extension baselines, focusing on open-sourced LLMs that underwent fine-tuning subsequent to positional interpolation techniques such as PI, NTK, and YaRN. The baselines encompass Together-32k~\cite{together}, Code LLaMA~\cite{codellama}, LongLoRA-full-FT-100k~\cite{longlora}, YaRN-LLaMA, and YaRN-Mistral~\cite{yarn}.

\subsection{Main Results}

\label{sec:mainresult}
\textbf{Language modeling on long sequences up to 256k tokens}. Our initial evaluation focuses on the performance of our approach against leading long-context large language models (LLMs) when tested on sequences extending to 256k tokens. To illustrate the robustness and versatility of our method, we utilize two datasets: the test splits from Proof-pile~\cite{pg19} and PG19~\cite{pile}. Perplexity is measured across different context sizes employing a sliding window mechanism of 256 tokens. For PG19, analysis is performed on the complete set of 100 test documents. In the case of Proof-pile, we adopt the procedure used in YaRN~\cite{yarn}, selecting 10 random samples, each with a minimum sequence length of 128k tokens.

Tables~\ref{tbl:proof-pile} and \ref{tbl:pg19} present perplexity results for LLaMA2 and Mistral, each augmented through various interpolation strategies, on the Proof-pile and PG19 benchmarks, respectively. Two main insights emerge: \textbf{(1)} the perplexity of our extended models consistently decreases as the evaluation length expands from 4k to 256k, indicating effective utilization of longer input sequences. \textbf{(2)} Despite operating with a context window that is 16$\times$ larger than the standard, which is generally associated with diminished performance on shorter sequences, our LongRoPE-2048k models consistently surpass leading baselines for context lengths up to 256k.

\textbf{Long sequence language modeling beyond 2000k}. To assess performance on exceptionally lengthy texts, we utilize the Books3~\cite{pile} dataset. For the sake of evaluation efficiency, we randomly sample 20 books, each with a length greater than 2048k, and apply a sliding window of 256k tokens.

Referencing Table~\ref{tbl:books3}, LongRoPE effectively expands the context capacity of both LLaMA2-7B and Mistral-7B models to 2048k tokens, while maintaining perplexity scores on par with, or better than, baseline models for shorter sequences in the 8k-128k range. Distinct performance characteristics are evident when comparing the 2048k variants of LLaMA2 and Mistral. At shorter sequence lengths, Mistral shows superior results relative to baselines; however, its perplexity rises above 7 when sequences surpass 256k tokens. LLaMA2 behaves in line with expectations: its perplexity consistently drops as the context length increases, with only minor upticks at 1024k and 2048k. Additionally, LongRoPE-2048k, when fine-tuned on LLaMA2 at a 256k length, yields better outcomes than at 128k, which can be attributed to the relatively smaller secondary extension ratio (specifically, 8$\times$ versus 16$\times$). In contrast, Mistral displays improved performance when the fine-tuning context window is set at 128k. This discrepancy arises mainly because Mistral's fine-tuning for 128k and 256k windows uses a 16k training length in accordance with YaRN's protocol, impacting its effectiveness in extending the context window further after fine-tuning.

\textbf{Passkey retrieval}. Here, we investigate how the size of the context window affects model performance on generation tasks. Following the methodology of~\cite{passkey}, we adopt a synthetic passkey retrieval task in which the model must locate a randomly assigned passkey (a five-digit numeral) embedded within an extensive text. Further specifics of the prompt template are provided in the appendix. The evaluation consists of 10 runs, where in each run the passkey is positioned at a randomly chosen location, uniformly sampled from the possible range of the context length used during evaluation.

Figure~\ref{fig:passkey} illustrates retrieval accuracy benchmarks against competing models. The accuracy of prior approaches declines sharply to zero once sequence length exceeds 128k tokens. In sharp contrast, even when tasked with retrieving a passkey amid millions of tokens, LongRoPE-LLaMA2-2048k (ft=256k) consistently achieves high retrieval rates ($\ge$90\%) across the full range from 4k to 2048k tokens. Meanwhile, LongRoPE-Mistral-2048k (ft=128k) maintains perfect accuracy through 1800k tokens, with a decrease to 60\% at 2048k, in agreement with observations from Table~\ref{tbl:books3}, where perplexity experiences a mild uptick at 2048k.

\textbf{Standard benchmarks within original context window}. LongRoPE-2048k models are assessed on their original context window through the Hugging Face Open LLM Leaderboard~\cite{huggingfaceboard}, considering both zero-shot and few-shot scenarios. Specifically, evaluation is conducted with 25-shot ARC-Challenge~\cite{allenai:arc}, 10-shot HellaSwag~\cite{zellers2019hellaswag}, 5-shot MMLU~\cite{mmlu}, and 0-shot TruthfulQA~\cite{lin2021truthfulqa}.

Table~\ref{tbl:commonsense} demonstrates that our models perform similarly to the original benchmark, which was optimized for a shorter context window, and surpass the original Mistral on TruthfulQA by +0.5\%. While LongRoPE-LLaMA2-2048k, fine-tuned with a 256k context, exhibits a modest decline in performance, it generally stays within acceptable limits across the majority of tasks.

\subsection{Ablation Results}

\textbf{Assessing the impact of the second positional interpolation.} Our progressive extension approach incorporates an additional, non-uniform positional interpolation by leveraging our search algorithm on LLMs that have been further fine-tuned. To demonstrate the utility of this method, we perform experiments on our fine-tuned LLaMA2-256k model and expand it to 512k, 1024k, and 2048k sequence lengths utilizing PI and YaRN. Results displayed in Table~\ref{tbl:secondsearch} indicate that our non-uniform positional interpolation maintains perplexity at a stable level. By comparison, applying PI and YaRN leads to a marked escalation in perplexity as the sequence extension increases.

\textbf{Short-context recovery effectiveness.} To address the decline in performance at reduced context lengths, we modify the RoPE scaling factors for LongRoPE-2048k using our search methodology. In this process, we intentionally lower the upper bound for scale factors to minimize interpolation at shorter lengths, specifically 4k and 8k. Table~\ref{tbl:shortperformance} provides a comparison of perplexity values for LongRoPE-LLaMA2-2048k evaluated on the Proof-pile dataset at both 4k and 8k, alongside averaged LLM benchmark accuracy scores. The outcomes demonstrate a notable boost in performance for shorter context lengths.

\textbf{Assessment of the two types of non-uniformities}. To understand the individual effect of each non-uniformity, we conduct ablation studies. We perform two sets of experiments: (i) expanding LLaMA2-7B to shorter sequence lengths of 16k and 32k using three distinct strategies—PI, optimizing solely the RoPE dimension, and optimizing both non-uniformities; (ii) scaling a fine-tuned LLaMA2 model trained on 256k-length sequences up to 2048k, employing a comparable methodology. Perplexity scores are measured prior to any fine-tuning. According to Table~\ref{tbl:searchdimension}, adjusting the RoPE dimension for non-uniformity yields a considerable decrease in perplexity relative to linear interpolation with PI. Introducing non-uniformity in token position results in noticeable improvements at 16k and 32k, while at 2048k, its effect diminishes—likely because of the sheer sequence length. At such scales, retaining only the early tokens without interpolation fails to offer meaningful benefit, and we highlight this as an area for further exploration.

\section{Related Works}

Beyond position interpolation methods, this section also covers research on alternative strategies.

\textbf{Retrieval-based} methods incorporate an external memory mechanism to retain extensive prior context, utilizing retrieval components to access pertinent documents during inference~\cite{longllama,wang2023augmenting,borgeaud2022improving}. Such techniques often require substantive changes to the underlying LLM architecture. By comparison, our approach remains streamlined, only involving minimal adjustments to positional embeddings. Additionally, our methodology extends to a broader set of tasks involving extended context, including long document summarization and few-shot learning, rather than being restricted to retrieval scenarios.

\textbf{Attention-based context window extensions}. In addition to positional embedding interpolation, another approach leverages modifications in attention mechanisms to extend the context window within the original LLM input constraints~\cite{han2023lminfinite, streamingllm,parallelwindow}. This strategy centers on controlling the increase in attention complexity that arises from introducing new token positions by employing innovative attention masks. Such techniques work in tandem with positional interpolation methods, serving as complementary solutions.

\textbf{Fine-tuning based approaches} concentrate on strategies for adapting pre-trained LLMs with altered position embeddings to support extended contexts. For example, studies such as Code LLaMA~\cite{codellama}, LLaMA2 Long~\cite{xiong2023effective}, and ScaledRoPE~\cite{liu2023scaling} implement fine-tuning by selecting a substantially increased RoPE base value and targeting specific context lengths during training. The approach introduced in this work accommodates a range of target lengths and is capable of handling sequences exceeding 2M tokens. Furthermore, since scaling fine-tuning to very large context windows (e.g., over 128k tokens) is computationally intensive, methods like LongLoRA~\cite{longlora} and PoSE~\cite{zhu2023pose} have been introduced to reduce the associated resource demands. Notably, the proposed method is complementary to these efficient fine-tuning techniques.

\section{Conclusion}

In this study, we introduce LongRoPE, a technique that significantly increases the context length of LLMs up to an extraordinary 2048k, all while preserving their performance with the original, shorter context range. Our approach leverages two distinct non-uniformities within RoPE positional embeddings, discovered via an efficient evolutionary search process. This strategy delivers two primary advantages: it serves as an excellent initialization for fine-tuning, and it facilitates an 8$\times$ extension of the context window without requiring fine-tuning. Furthermore, we develop a progressive extension method, utilizing LLMs fine-tuned at 256k context length to scale up to a 2048k window size, foregoing additional fine-tuning. Through comprehensive experimentation, we demonstrate LongRoPE's efficacy. We anticipate that LongRoPE-2048k will enable numerous applications requiring extended contexts and stimulate further investigation in this domain.

\section*{Broader Impacts} This study aims to contribute to the progress of Machine Learning as a discipline. While our research may have various implications for society, we do not believe that any particular issues require special emphasis at this time.

	\balance

\end{document}